\documentclass[11pt,a4paper]{article}
\usepackage[T1]{fontenc}
\usepackage[utf8]{inputenc}
\usepackage[english]{babel}
\usepackage{amsmath,amssymb,amsthm}
\usepackage[margin=2.7cm]{geometry}

\newtheorem{theorem}{Theorem}[section]
\newtheorem{proposition}[theorem]{Proposition}
\newtheorem{lemma}[theorem]{Lemma}
\newtheorem{corollary}[theorem]{Corollary}
\theoremstyle{definition}
\newtheorem{definition}[theorem]{Definition}
\theoremstyle{remark}
\newtheorem{remark}[theorem]{Remark}
\newtheorem{example}[theorem]{Example}

\newcommand{\Dinj}{\Delta_{\mathrm{inj}}}
\newcommand{\Tr}{\mathrm{Tr}}
\newcommand{\supp}{\mathrm{supp}}
\newcommand{\rk}{\mathrm{rk}}
\newcommand{\sel}{\mathrm{sel}}
\newcommand{\im}{\mathrm{im}}
\newcommand{\Acc}{\mathrm{Acc}}
\newcommand{\id}{\mathrm{id}}
\newcommand{\N}{\mathbb{N}}

\title{A constructive theory of dilators}
\author{}
\date{}

\begin{document}
\maketitle

\begin{abstract}
We present a constructive version of part of Girard's theory of dilators
(\emph{Proof Theory and Logical Complexity}, chapters 8 and 9): the normal
form, the extension to linear orders, the well-foundedness of the dilators
built by sum, product and exponentiation in base $2$, the link with the
slow growing and Hardy hierarchies up to the Howard ordinal, morphisms,
connected components and kinds. All objects are finite or decidable,
well-foundedness is expressed by accessibility, and no proof uses the
excluded middle. Everything has been checked in type theory, without
axioms (function extensionality is used, as a hypothesis, only for two
statements about trees of level $\Omega$).
\end{abstract}

\section{Introduction}

Girard presents the theory of dilators in a classical setting: the
category of ordinals, functors preserving direct limits and pull-backs,
well-orders. Many of the arguments are nevertheless finite and
combinatorial, and Girard's normal form (a point of $F(x)$ is written
uniquely as $(z;x_0,\dots,x_{n-1};x)_F$) suggests that the whole theory
reduces to finite, decidable data. The aim of this note is to make this
explicit.

The following choices are made throughout.
\begin{itemize}
\item We work on the category $\Dinj$ of finite ordinals with strictly
  increasing maps, coded by words. A pre-dilator is a functor
  $\Dinj \to \Dinj$ preserving pull-backs. It is thus given by finite data
  in each degree, and every property of it in a fixed degree is
  decidable.
\item A linear order is a type with a decidable comparison with values in
  $\{<,=,>\}$. The extension of a pre-dilator to such an order is again a
  decidable linear order.
\item Well-foundedness is accessibility: $x$ is accessible if all its
  predecessors are, as an inductive definition. An order is well-founded
  if all its points are accessible. The absence of infinite descending
  sequences is never used.
\item A disjunction is asserted only when we can tell which case holds.
  Where Girard concludes ``$F' = 0$ or $F'' = 0$'' with no possible
  decision, we state the negative form.
\end{itemize}
We point out along the way where the constructive version differs from
the book.

\section{Strictly increasing maps as words}

Write $[n] = \{0, \dots, n-1\}$. A strictly increasing map
$f : [m] \to [n]$ is determined by its image, hence by the word
$w \in \{o,i\}^n$ whose letter $y$ is $i$ exactly when $y \in \im f$.
Conversely every word $w$ of length $|w| = n$ with $\#w = m$ letters $i$
is a strictly increasing map $[m] \to [n]$. On words one defines by
recursion:
\begin{itemize}
\item the composite $v \circ w$ (when $\#v = |w|$): replace the $k$-th
  letter $i$ of $v$ by the $k$-th letter of $w$;
\item the intersection $v \wedge w$ of two words of the same length,
  letter by letter; it is the pull-back of $v$ and $w$ in $\Dinj$;
\item the rank $\rk_w(y)$, the number of letters $i$ before position $y$
  (the preimage of $y$), and the selection $\sel_w(z)$, the position of
  the $z$-th letter $i$ (the map $w$ itself);
\item the inclusion $v \le w$ of images.
\end{itemize}
Everything is decidable. The basic lemmas are: $\sel_w$ is strictly
increasing on $[\#w]$, $\sel_{v \circ w} = \sel_v \circ \sel_w$,
$\sel_w(\rk_w(y)) = y$ if $y \in \im w$, two words of the same length and
the same image are equal, and a word $w$ such that $s \circ w$ has the
same image as $s$ is the identity.

\section{Pre-dilators and normal form}

\begin{definition}
A \emph{pre-dilator} is given by $D_0 : \N \to \N$ and a map $D_1$ from
words to words such that $|D_1 w| = D_0 |w|$, $\#D_1 w = D_0(\#w)$,
$D_1(\id_n) = \id_{D_0 n}$, $D_1(v \circ w) = D_1 v \circ D_1 w$ and
$D_1(v \wedge w) = D_1 v \wedge D_1 w$.
\end{definition}

\begin{remark}
Girard's Definition 8.2.16 also requires monotonicity: $f \le g$ implies
$F(f) \le F(g)$. We do not need it. For a dilator it follows from
well-foundedness (8.2.15), by an argument on $\omega^{2y}$; everything
below holds without it.
\end{remark}

Let $x \in D(p) = [D_0 p]$. A \emph{support} of $x$ is a word
$u : [k] \to [p]$ with $x \in \im D_1 u$.

\begin{theorem}[normal form, Girard 8.2.3]\label{th:nf}
\begin{enumerate}
\item The supports of $x$ are closed under intersection, and $x$ has a
  least support $\supp(x)$, which is computable.
\item The point $z = \rk_{D_1(\supp x)}(x)$ of $D(n)$, $n = \#\supp x$, is
  in the \emph{trace} $\Tr(n)$: it is in the image of no $D_1 v$ with
  $v : [k] \to [n]$, $v \neq \id$.
\item If $x = \sel_{D_1 w}(z')$ with $z' \in \Tr$, then $w = \supp(x)$ and
  $z' = z$.
\end{enumerate}
\end{theorem}

\begin{proof}
(1) If $x \in \im D_1 u \cap \im D_1 v$, then $x \in \im(D_1 u \wedge D_1
v) = \im D_1(u \wedge v)$: this is the only use of pull-backs. To compute
the least support, let $g(u) = u$ if $u$ is a support and $g(u) = \id_p$
otherwise (this is decidable), and
\[
\supp(x) = \bigwedge_{u \in \{o,i\}^p} g(u).
\]
The intersection ranges over a \emph{finite} set of words; it is a support
by (1), and it is contained in every support $u$ since $g(u) = u$.
(2) If $z = \sel_{D_1 v}(z'')$, then $\supp(x) \circ v$ is again a support
of $x$, hence contains $\supp(x)$, and $v = \id$. (3) The word $w$ is a
support, so $\supp(x) = w \circ r$ for some word $r$; and
$\rk_{D_1 w}(x) \in \im D_1 r$ forces $r = \id$ by the hypothesis on $z'$.
\end{proof}

\begin{remark}
In Girard, $\supp(x)$ is the intersection of all supports in $ON$, an
infinite family. On $\Dinj$ it is indexed by the $2^p$ words of length
$p$, and the normal form becomes an algorithm. In particular the trace is
decidable: $z \in \Tr(n)$ if and only if $\supp(z) = \id_n$.
\end{remark}

\begin{corollary}[denotation form]
The maps $x \mapsto (\supp x, z)$ and $(w, z) \mapsto \sel_{D_1 w}(z)$ are
inverse bijections
\[
D(p) \;\cong\; \sum_{n \le p} \Tr(n) \times \Dinj([n],[p]).
\]
\end{corollary}

\section{Extension to decidable linear orders}

Let $X$ be a decidable linear order. A strictly increasing map
$[n] \to X$ is a sorted list of length $n$.

\begin{definition}
$D(X)$ is the set of \emph{denotations} $(a, z)$, where $a$ is a sorted
list of $X$ and $z \in \Tr(|a|)$. To compare $(a, z)$ and $(a', z')$, form
the merge $c$ of $a$ and $a'$, with $a = c \circ u$ and $a' = c \circ u'$,
and compare $\sel_{D_1 u}(z)$ and $\sel_{D_1 u'}(z')$ in $\N$.
\end{definition}

\begin{theorem}
\begin{enumerate}
\item The result does not depend on the common refinement $c$.
\item The comparison makes $D(X)$ a decidable linear order.
\item A strictly increasing $f : X \to Y$ induces
  $D(f)(a, z) = (f \circ a, z)$, strictly increasing; $D$ is a functor,
  and the image of $D(f)$ consists of the denotations whose support lies
  in the image of $f$ (preservation of pull-backs).
\item Every denotation is the image of a point of $D(n)$ along its own
  list (direct limits). For $X = \omega$, the denotations bounded by $k$
  form a copy of $D(k)$.
\end{enumerate}
\end{theorem}

\begin{proof}[Sketch]
Two common refinements factor through their merge, and a sorted list is
injective; one concludes by $\sel_{D_1(ev)} = \sel_{D_1 e} \circ
\sel_{D_1 v}$ and the strict monotonicity of $\sel_{D_1 e}$. Transitivity
is read in a common refinement of the three lists. If two denotations give
the same point of $D(|c|)$, they are equal by the denotation form.
\end{proof}

\section{Well-foundedness}

\begin{definition}
For a decidable relation $<$ on $A$, accessibility $\Acc$ is the least
predicate such that $x$ is accessible as soon as all $y < x$ are. $A$ is
\emph{well-founded} if every point is accessible. A pre-dilator $D$ is a
\emph{dilator} if $D(X)$ is well-founded for every well-founded decidable
linear order $X$.
\end{definition}

This definition quantifies over all orders $X$ in a universe; it is a
statement of a higher universe. The induction principle associated with
$\Acc$ is the only form of transfinite induction used.

\begin{lemma}[inverse image]
If $\varphi : A \to B$ preserves $<$ and $B$ is well-founded, so is $A$.
The ordered sum and the lexicographic product of two well-founded orders
are well-founded, and so is $\N$.
\end{lemma}

\begin{theorem}\label{th:poly}
$\mathrm{Id}$, $1$, $D + E$, $D \cdot E$ are dilators if $D$ and $E$ are.
\end{theorem}

\begin{proof}
In each case $D(X)$ is mapped into a known order, preserving the order.
The trace of $\mathrm{Id}$ is $\{0\} = \Tr(1)$, so $\mathrm{Id}(X) \cong
X$. For the sum, $(D + E)(n) = D(n) + E(n)$ and the trace of $D + E$ is
that of $D$ followed by that of $E$, shifted; $(D+E)(X)$ maps into the
ordered sum $D(X) + E(X)$. For the product, $(D \cdot E)(n)$ is $E(n)$
copies of $D(n)$, ordered by the Kronecker product of words. A point
$z = e \cdot D_0(n) + j$ (Euclidean division) gives the points $e$ of
$E(n)$ and $j$ of $D(n)$, which are pushed into $E(X)$ and $D(X)$; the
order is lexicographic, $E$ first.
\end{proof}

\section{Exponentiation}

\subsection{$2^X$ is well-founded}

For a linear order $Y$, $2^Y$ is the set of finite subsets, compared by
the largest element of the symmetric difference. It is represented by the
strictly decreasing lists, compared lexicographically (the first element
is the largest, and a proper prefix is smaller).

\begin{theorem}\label{th:2X}
If $Y$ is well-founded, so is $2^Y$.
\end{theorem}

\begin{proof}
The empty list is the least element, hence accessible. Let $P(x)$ be:
every decreasing list with head $x$ is accessible. We prove $P(x)$ by
induction on the accessibility of $x$. Assume $P(y)$ for all $y < x$.
\begin{itemize}
\item Every decreasing list whose head is $< x$ is accessible: it is the
  empty list, or a list with head $y < x$.
\item Let $l$ be a decreasing list with elements $< x$; it is accessible
  by the previous point. We show, by induction on the accessibility of
  $l$, that $x :: l$ is accessible. A predecessor of $x :: l$ is the empty
  list, or $y :: k$ with $y < x$ (accessible by $P(y)$), or $x :: k$ with
  $k < l$ (accessible by the induction hypothesis on $l$).
\end{itemize}
Since every decreasing list is empty or of the form $x :: l$, all of them
are accessible.
\end{proof}

\begin{remark}
Girard recalls (§9.6) that the statement ``$X$ well-ordered $\Rightarrow$
$2^X$ well-ordered'' is equivalent, over a weak theory, to
$(\Sigma^0_1\text{-CA})^*$. With accessibility as an inductive definition,
the proof is an elementary double induction: the logical strength lies in
the existence of inductive definitions, not in a comprehension principle.
\end{remark}

\subsection{The pre-dilator $2^D$}

On finite ordinals, $2^n$ is the number $2^n$, a subset $s$ of $[n]$ being
the number $\sum_{k \in s} 2^k$. The order of numbers is then the order of
$2^{[n]}$. For $w : [m] \to [n]$, $2^w : [2^m] \to [2^n]$ sends the subset
$s$ to its image $w(s)$; on words, position $j$ is in the image of $2^w$
if and only if the subset $j$ is contained in $\im w$. The first letter of
$w$ corresponds to the least significant bit, which gives the recursion
\[
2^{\varnothing} = i, \qquad 2^{o\,w} = \mathrm{lo}(2^w), \qquad
2^{i\,w} = \mathrm{dbl}(2^w),
\]
where $\mathrm{lo}$ puts an $o$ after each letter and $\mathrm{dbl}$
doubles each letter. Composition and intersection commute with
$\mathrm{lo}$ and $\mathrm{dbl}$, so $2^-$ is a pre-dilator, and so is
$2^D = 2^- \circ D$.

\begin{theorem}\label{th:2D}
If $D$ is a dilator, so is $2^D$.
\end{theorem}

\begin{proof}
A denotation $(a, z)$ of $(2^D)(X)$ is a subset $s$ of $D(|a|)$ (the bits
of $z$). We associate with it the decreasing list $\varphi(a,z)$ of the
points $D(a)(t)$ of $D(X)$, for $t \in s$. By Theorem~\ref{th:2X} and the
inverse image lemma, it suffices to show that $\varphi$ preserves the
order. Let $c$ be the merge of $a$ and $a'$, with $a = c \circ u$. Three
combinatorial facts suffice:
\begin{enumerate}
\item the point of $2^{D(|c|)}$ image of $s$ is the subset $D_1 u(s)$,
  that is $\sel_{2^W}(\mathrm{num}\,s) = \mathrm{num}(W \circ s)$ for a
  word $W$;
\item comparing two numerals of the same length amounts to comparing
  lexicographically the decreasing lists of their bits equal to $1$;
\item the bits equal to $1$ of $W \circ s$ are the images by $\sel_W$ of
  those of $s$, in the same order.
\end{enumerate}
Finally, two images $\sel_{D_1 u}(t)$ and $\sel_{D_1 u'}(t')$ compare as
the points $D(a)(t)$ and $D(a')(t')$ of $D(X)$.
\end{proof}

\begin{corollary}
The pre-dilators built from $\mathrm{Id}$ and $1$ by $+$, $\cdot$ and
$2^-$ are dilators; in particular $D(\omega)$ is well-founded.
\end{corollary}

\section{Trees and hierarchies}

We use the Brouwer trees $\mathcal O$ ($0$, successor, limit of a
sequence) and $\mathcal O_2$ (in addition, limit of a family indexed by
$\mathcal O$), with the slow growing hierarchy $G_1 : \mathcal O \to \N
\to \N$, the Hardy hierarchy $H_1$, and their analogues $G_2 : \mathcal
O_2 \to \N \to \mathcal O$, $H_2 : \mathcal O_2 \to \mathcal O \to
\mathcal O$.

With each polynomial dilator $p$ we associate a tree $t(p)$:
$\mathrm{Id} \mapsto \omega$, $1 \mapsto 1$, sum and product of trees, and
$2^b$ with $2^{b+1} = 2^b + 2^b$. We have
\[
G_1(t(p))(k) = |D_p(k)|,
\]
which is Girard's statement 9.A for these dilators: the slow growing
hierarchy along $D(\omega)$ gives back $D$. The tree $t_2(p) \in \mathcal
O_2$ (where $\mathrm{Id} \mapsto \Omega$) satisfies the condition on
fundamental sequences needed for Wainer's collapsing theorem, and
$G_2(t_2(p))(n) = t(p)$.

\begin{theorem}
Let $T_0 = \mathrm{Id} + 1$ and $T_{k+1} = 2^{T_k}$. The limits of the
trees are $\varepsilon_0 = \sup_k t(T_k)$ and $\varepsilon_{\Omega+1} =
\sup_k t_2(T_k)$; let $\eta_0 = H_2(\varepsilon_{\Omega+1})(\omega)$ (the
Howard ordinal, in Girard's sense, 9.A). Then, for all $n$,
\[
G_1(\eta_0)(n) = H_1(\varepsilon_0)(n).
\]
\end{theorem}

Function extensionality is used in one place only: the commutation of
$G_2$ with the arithmetic under the limits indexed by $\mathcal O$. It is
taken as a hypothesis.

\section{Morphisms}

\begin{definition}
A natural transformation $T : D \to E$ is a family of words
$t_n : [D_0 n] \to [E_0 n]$ such that $E_1 w \circ t_m = t_n \circ D_1 w$
for $w : [m] \to [n]$.
\end{definition}

\begin{proposition}[Girard 8.2.5 (ii)]
If $z \in \Tr_D(m)$, then $T(m)(z) \in \Tr_E(m)$, and $T$ sends the
denotation $(w, z)$ to $(w, T(m)(z))$.
\end{proposition}

\begin{proof}
Let $v : [k] \to [m]$ with $T(m)(z) \in \im E_1 v$. The equality $v = \id$
is decidable; assume $v \neq \id$. Then $v$ misses a point $p$, and the two
injections $g_1, g_2 : [m] \to [m+1]$ that skip $p$ and $p+1$ are distinct
and agree on $v$. So $E(g_1)$ and $E(g_2)$ agree on $T(m)(z)$, and by
naturality and injectivity of $T(m+1)$, $D(g_1)(z) = D(g_2)(z)$. Since $z$
is in the trace, $g_1$ and $g_2$ are both the least support of this point
(Theorem~\ref{th:nf}), a contradiction. Hence $\neg\neg(v = \id)$, and
$v = \id$ by decidability.
\end{proof}

\begin{proposition}[Girard 8.3.5 (ii)]
Two natural transformations $D \to E$ that agree on the traces are equal.
\end{proposition}

\begin{proof}
Every $x \in D(n)$ is $\sel_{D_1 w}(z)$ with $z$ in the trace, and
$T(n)(x) = \sel_{E_1 w}(T(m)(z))$. Two words of the same length with the
same values under $\sel$ are equal.
\end{proof}

The inclusions of $D$ and of $E$ into $D+E$ are natural; their traces are
$z \mapsto z$ and $y \mapsto D_0(m) + y$.

\section{Connected components}

\begin{definition}
Two trace points $(m, z)$ and $(m', z')$ of $D$ are \emph{separated} if,
for all words $u, u'$ of the same length with $\#u = m$ and $\#u' = m'$,
we have $\sel_{D_1 u}(z) < \sel_{D_1 u'}(z')$. They are \emph{linked} if
they are separated in neither order.
\end{definition}

Being separated means that every denotation with trace $(m,z)$ is below
every denotation with trace $(m',z')$, in $D(X)$ for every $X$, whatever
the coefficients: this is $\S_D(m,z;m',z') = (0,+1)$ in Girard (8.4.18).

\begin{definition}
On $\Dinj$ an increasing bijection of $[n]$ is the identity, so two
isomorphic pre-dilators are equal. We write $D = F + G$ if $D_0 n = F_0 n
+ G_0 n$ and $D_1 w = F_1 w \cdot G_1 w$ (concatenation of words) for all
$n$ and $w$. $D$ is \emph{non-zero} if there are $n$ and a point of
$D(n)$. $D$ is \emph{connected} if it is non-zero and, for every
decomposition $D = F + G$, $F$ and $G$ are not both non-zero.
\end{definition}

\begin{remark}
Girard requires ``$F = 0$ or $G = 0$''. The statement $F = 0$, that is
$\forall n\; F_0 n = 0$, is not decidable, and a pre-dilator that is zero
in degree $n$ may be non-zero in a higher degree. The negative form is the
right constructive version; classically it is equivalent.
\end{remark}

\begin{theorem}[Girard 9.1.7]\label{th:917}
If $F$ and $G$ are non-zero, $F + G$ has two separated trace points.
Consequently a non-zero pre-dilator whose trace points are pairwise linked
is connected.
\end{theorem}

\begin{proof}
A point of $F(n)$ has a normal form with trace $(m, z)$, a point of
$G(n')$ a normal form with trace $(m', y)$. In $F + G$, $(m, z)$ and
$(m', F_0 m' + y)$ are trace points (images under the inclusions), and for
$u, u'$ of the same length $L$,
\[
\sel_{(F+G)_1 u}(z) = \sel_{F_1 u}(z) < F_0 L \le F_0 L + \sel_{G_1
u'}(y) = \sel_{(F+G)_1 u'}(F_0 m' + y). \qedhere
\]
\end{proof}

\begin{example}
$\mathrm{Id}$ and $1$ each have a single trace point, $(1,0)$ and
$(0,0)$; they are connected. $\mathrm{Id} + 1$ is not.
\end{example}

The construction of the components rests on the restriction to a set of
trace points.

\begin{theorem}[Girard 8.3.6]\label{th:836}
Let $c$ be a decidable set of trace points of $D$. Let $F(n) = \{x \in
D(n) \mid$ the trace of $x$ is in $c\}$ and $F(w)$ the restriction of
$D(w)$. Then $F$ is a pre-dilator, the inclusion $F \to D$ is a natural
transformation, and its trace is contained in $c$.
\end{theorem}

\begin{proof}
Let $\chi_n : [F_0 n] \to [D_0 n]$ be the inclusion. The key point is that
$D(w)$ preserves traces: if $x = D(s)(z)$ with $z$ in the trace, then
$D(w)(x) = D(w \circ s)(z)$, and by uniqueness of the normal form the
trace of $D(w)(x)$ is $(\#s, z)$. It follows that $F$ is a pull-back of
$D$:
\[
D_1 w \circ \chi_m = D_1 w \wedge \chi_n .
\]
We define $F_1 w$ by $\chi_n \circ F_1 w = D_1 w \circ \chi_m$. The
functorial identities are checked after composing on the left with
$\chi_n$, which can be cancelled, using for intersections that
composition distributes over $\wedge$.
\end{proof}

\begin{theorem}[cut]\label{th:cut}
Let $c$ be as above, such that every trace point in $c$ is separated from
(and below) every trace point outside $c$. Then $D = D|_c + D|_{\neg c}$.
In particular a connected pre-dilator has no proper cut.
\end{theorem}

\begin{proof}
If $x, y \in D(n)$ have their traces in $c$ and outside $c$ respectively,
then $x < y$: apply separation to the supports of $x$ and $y$. The points
of $D|_c(n)$ thus form an initial segment of $D(n)$, so that $\chi_n =
\id_{a} \cdot 0^{b}$ and $\chi'_n = 0^{a} \cdot \id_{b}$. The pull-back
property then identifies the two blocks of $D_1 w$ with $F_1 w$ and
$G_1 w$.
\end{proof}

Girard's Theorem 9.1.6 (unique decomposition of a dilator as a sum of
connected components indexed by an ordinal) is obtained by cutting along
the classes of the relation ``linked''; Theorems~\ref{th:836}
and~\ref{th:cut} are its constructive building blocks. What remains is the
transitivity of this relation (the ultrametric inequality 8.4.21), the
well-ordering of the quotient, and the indexed sum.

\section{Kinds}

Girard classifies dilators into four kinds (9.1.10): $0$; $1$ ($F = F' +
1$); $\omega$ (a sum indexed by a limit ordinal); $\Omega$ ($F = F' + F''$
with $F''$ connected and different from $1$). For a tree in $\mathcal
O_2$, the kind is the head constructor: $0$, successor, limit over $\N$,
limit over $\mathcal O$.

\begin{theorem}
Let $p$ be a polynomial dilator (with $2^-$). The tree $t_2(p)$ is of kind
$1$ or $\Omega$, never $0$ nor $\omega$. It is of kind $1$ if and only if,
for some (equivalently every) $n$, the largest point of $D_p(n+1)$ is a
constant, that is, a point coming from $D_p(0)$.
\end{theorem}

For every $F$, the largest point of $(F + 1)(n)$ is a constant; the
theorem thus identifies, for polynomial dilators, kind $1$ of the tree
with Girard's kind $1$. The kind of $2^D$ is $1$ exactly when $D$ is
constant.

\section{Finite and ideal elements}

In constructive domain theory one works with the finite (compact)
elements, and expresses results in terms of ideal elements, which are
determined by their finite approximations. The theory of dilators has
exactly this structure, and historically the two are related: Girard's
stable domains and coherence spaces were obtained by simplifying
dilators.

\medskip
\begin{tabular}{ll}
\emph{domain theory} & \emph{dilators} \\ \hline
compact elements & finite ordinals, $\Dinj$, words \\
information system & the data $(D_0, D_1)$ on words \\
ideal elements & $D(X)$, $X$ any decidable linear order \\
continuity & preservation of direct limits \\
stability (Berry) & preservation of pull-backs \\
minimal approximation & least support \\
trace of a stable map & trace $\{(n, z) \mid z \in \Tr(n)\}$ \\
points of the completion & denotations $(z; x_0 < \dots < x_{n-1}; X)$
\end{tabular}

\medskip
Every linear order is the directed union of its finite suborders, and a
pre-dilator, which preserves direct limits, is determined by its
restriction to $\Dinj$. This is the analogue of a continuous map being
determined by its action on compact elements.

\paragraph{Stability gives canonical names.}
In a Scott domain an ideal element is a class of approximations, with no
canonical representative. Here, by Theorem~\ref{th:nf}, every point of
$D(X)$ has a \emph{least} finite approximation, its support. So an ideal
point has a canonical name, the pair (support, trace point), and $D(X)$
can be defined directly by denotations, with decidable order and
equality, instead of as a quotient. Preservation of pull-backs plays the
role of Berry's stability condition, and it is what makes the ideal level
as computable as the finite one. Conceptually $D(X)$ is the colimit of the
$D(n)$ over the finite subsets of $X$ (a left Kan extension along
$\Dinj \to$ linear orders); the denotations compute this colimit.

\paragraph{Proofs on finite elements, statements on ideal elements.}
All the proofs above are about finite objects: words, $D(n)$, the merge
of two supports. Many of the statements are about ideal ones: $D(X)$ is a
decidable linear order, functorial in $X$, well-founded when $X$ is, and
the hierarchy theorems are about $D(\omega)$ and $D(\Omega)$. The results
on morphisms say that a natural transformation is determined by its action
on traces (Proposition 8.3.5 (ii)), as an approximable mapping is
determined on the basis. Even $\Omega$ is handled as an ideal ordinal: the
tree of $D(\Omega)$ has limits indexed by all countable trees, and the
collapsing theorem only uses it through its approximations.

\paragraph{Uniqueness of the ideal object.}
Call an \emph{extension} of $D$ to $X$ an ordered type $G$ with, for each
finite sorted subset $c$ of $X$, a map $\iota_c : D(|c|) \to G$ such that
$\iota_c \circ D(w) = \iota_{c \circ w}$, $\iota_c$ is strictly increasing,
and every point of $G$ is some $\iota_c(x)$. (Preservation of pull-backs
is not needed: it is already in $D$.)

\begin{theorem}[Girard 8.2.6, univalent form]\label{th:uniq}
Assume function extensionality and univalence. Every extension $G$ of $D$
to $X$ is equal, as an ordered type, to $D(X)$.
\end{theorem}

\begin{proof}
The map $\varphi(a, z) = \iota_a(z)$ preserves the comparison: one
computes in the merge $c$ of two supports, using naturality and the
monotonicity of $\iota_c$. It is injective because $D(X)$ is linear, and
surjective because $\iota_c(x) = \varphi(c \circ w, z)$, where $(w, z)$ is
the normal form of $x$ in $D(|c|)$. An equivalence preserving the
comparison is an equality of ordered types, by univalence and the
structure identity principle.
\end{proof}

Univalence and function extensionality are hypotheses (types), not
axioms, and nothing at the finite level depends on them. So the three
levels are separated: the finite level is decidable, the ideal level is
described by its finite approximations, and the univalent level turns
``determined by'' into equality.

\paragraph{Where Girard's text looks non-constructive.}
These are the places where an ideal element is treated as a completed
whole: deciding the kind of a dilator, deciding $F = 0$, the cofinality of
$F(\mathrm{On})$. The remedy is the one of domain theory: replace the
global decision by a statement about approximations. Connectedness,
stated negatively, is an example.

\section{Summary}

The above shows that a substantial part of Girard's theory is
constructive, and often more simply so than in the book:
\begin{itemize}
\item the normal form is an algorithm, because the intersection of
  supports ranges over a finite family;
\item the extension to linear orders and all comparisons are decidable;
\item well-foundedness, expressed by accessibility, propagates by
  elementary inductions, including for $2^X$;
\item Girard's proofs by contradiction (8.2.5) concern decidable
  properties and can be eliminated;
\item connectedness must be stated in negative form, and it is the only
  notion of chapter 9 treated here that is not decidable.
\end{itemize}
Not treated: 9.1.6 in full, kinds for an arbitrary dilator, the order type
of $t(p)$ as the ordinal of $D_p(\omega)$, and beyond, $\Pi^1_2$-
completeness, the well-ordering of dilators and the functor $\Lambda$.

\appendix
\section{Correspondence with the formalization}

The file \texttt{dilator.seg} (about 1350 declarations) is checked by a
type-theory checker in which the only recursion is the elimination of
inductive types, with no axioms.

\medskip
\begin{tabular}{ll}
Theorem~\ref{th:nf} and corollary & \texttt{PreDil.Point}, \texttt{fwd}, \texttt{bwd} \\
Extension to $X$ & \texttt{LinOrd.DenOrd}, \texttt{Fun} \\
Theorem~\ref{th:poly} & \texttt{IdWF}, \texttt{OneWF}, \texttt{SumWF}, \texttt{ProdWF}, \texttt{polDil} \\
Theorem~\ref{th:2X} & \texttt{Exp2.wfE} \\
Theorem~\ref{th:2D} & \texttt{ExpDil}, \texttt{ExpWF} \\
Trees, $\eta_0$ & \texttt{hier}, \texttt{FE.hierH}, \texttt{FE.hierEta} \\
Morphisms & \texttt{Nat.trT}, \texttt{Nat.denT}, \texttt{NatEq.Pts.eqT} \\
Theorem~\ref{th:917} & \texttt{ConnSum}, \texttt{ConnCrit.connected} \\
Theorem~\ref{th:836} & \texttt{Restr} \\
Theorem~\ref{th:cut} & \texttt{SplitCut.isSum}, \texttt{noCut} \\
Kinds & \texttt{treeK}, \texttt{kindPol}, \texttt{kind1} \\
Theorem~\ref{th:uniq} & \texttt{UA.sip}, \texttt{UA.Ext.Uniq.uniq}
\end{tabular}

\end{document}
